\begin{figure}[H]\centering
\begin{tikzpicture}[x=1cm,y=1cm]
\begin{scope}
\draw[figblue,thick] (0,0)--(5,2);
\draw[figred,thick] (0,-0.25)--(5,2.15);
\fill[figgray] (3.125,1.25) circle (2pt);
\fill[figamber] (0.9,0.3) circle (2.5pt);
\node[above left,font=\small] at (3.125,1.25) {$\mathbf x$};
\node[below,font=\small] at (0.9,0.3) {$\overline{\mathbf x}$};
\node[align=center,font=\small] at (2.5,-1) {(a) Cerca de ambas rectas,\\lejos de la intersección};
\end{scope}
\begin{scope}[xshift=6.4cm]
\draw[figblue,thick] (0,0)--(4,2);
\draw[figred,thick] (1,3)--(3,-1);
\fill[figgray] (2,1) circle (2pt);
\fill[figamber] (2.2,1.05) circle (2.5pt);
\node[above left,font=\small] at (2,1) {$\mathbf x$};
\node[right,font=\small] at (2.2,1.05) {$\overline{\mathbf x}$};
\node[align=center,font=\small] at (2,-1.5) {(b) Direcciones separadas:\\cerca de ambas y de la solución};
\end{scope}
\end{tikzpicture}
\caption{Comparación cualitativa dibujada por el docente. Pendientes y separaciones se eligen para hacer visible el mecanismo; no son los coeficientes del ejemplo numérico, que no se recuperan del audio. La cercanía geométrica ilustra el incumplimiento pequeño de ambas ecuaciones.}
\end{figure}
